\backmatter
\AfterBibliographyPreamble{%
Every measurement in this work is built on somebody's published result, and the
first purpose of this chapter is that those people are credited by name. A surname
in a footnote is not a citation. The second purpose is that the credit is checkable:
each source is registered in \texttt{anchor\_sift\_citations}, where the copy sits
itself, and this bibliography is generated from that registry and pinned to one
commit of it. A reference therefore resolves not to a title but to exact bytes.

Credit here does not depend on agreement. Where this work reproduces a published
result it says so, and where it fails to reproduce one it says that too, under the
same names, because a refutation rests on the original as much as a confirmation does.

\begin{description}
\item[Registry commit] \texttt{1ecbfcbce118b267de23756a6037a9978e6d855a}
\item[Dated] 2026-09-25
\item[Inventory] \texttt{MANIFEST.tsv}, signed; the hashes below are its rows
\end{description}

A reader who disagrees with a number in this research paper can ask which copy it was checked
against, and the hash answers exactly. Two copies of a paper are rarely the same
bytes.

These results are used in the work and their authors are not yet properly
credited: the registry holds a key and no bibliographic fields.
\par\medskip\raggedright}
\KOMAoptions{bibliography=totoc}
\begin{thebibliography}{14}

\bibitem{src:Bombieri}
E. Bombieri. \emph{Problems of the Millennium: the Riemann Hypothesis}. Clay Mathematics Institute, Millennium Prize Problem statement, \texttt{https:/\allowbreak{}/\allowbreak{}www.claymath.org/\allowbreak{}wp-content/\allowbreak{}uploads/\allowbreak{}2022/\allowbreak{}05/\allowbreak{}riemann.pdf}.\newline
{\footnotesize Copy held \texttt{sources/\allowbreak{}mathematics/\allowbreak{}millennium\_\allowbreak{}problems/\allowbreak{}clay\_\allowbreak{}bombieri\_\allowbreak{}riemann\_\allowbreak{}hypothesis.pdf}, SHA-256 \texttt{1454b2909f99271726ffb68b056aef45b7d3e6893a66282cad596339d69bafa9}.}

\bibitem{src:Clay-Mathematics-Institute}
Clay Mathematics Institute. \texttt{https:/\allowbreak{}/\allowbreak{}www.claymath.org/\allowbreak{}millennium-problems/\allowbreak{}}.

\bibitem{src:Cook}
Stephen Cook. \emph{The P versus NP Problem}. Clay Mathematics Institute, Millennium Prize Problem statement, \texttt{https:/\allowbreak{}/\allowbreak{}www.claymath.org/\allowbreak{}wp-content/\allowbreak{}uploads/\allowbreak{}2022/\allowbreak{}06/\allowbreak{}pvsnp.pdf}.\newline
{\footnotesize Copy held \texttt{sources/\allowbreak{}mathematics/\allowbreak{}millennium\_\allowbreak{}problems/\allowbreak{}clay\_\allowbreak{}cook\_\allowbreak{}p\_\allowbreak{}versus\_\allowbreak{}np.pdf}, SHA-256 \texttt{018f3d473d16c35e807e8cdfbd0bed3ccf7a56452e3e10c6b1d84b56eaf2cf59}.}

\bibitem{src:Deligne}
Pierre Deligne. \emph{The Hodge Conjecture}. Clay Mathematics Institute, Millennium Prize Problem statement, \texttt{https:/\allowbreak{}/\allowbreak{}www.claymath.org/\allowbreak{}wp-content/\allowbreak{}uploads/\allowbreak{}2022/\allowbreak{}06/\allowbreak{}hodge.pdf}.\newline
{\footnotesize Copy held \texttt{sources/\allowbreak{}mathematics/\allowbreak{}millennium\_\allowbreak{}problems/\allowbreak{}clay\_\allowbreak{}deligne\_\allowbreak{}hodge\_\allowbreak{}conjecture.pdf}, SHA-256 \texttt{e308d945ea3cf5dad8b187a06509013712c467b589039eb41b365cc4c988f0c8}.}

\bibitem{src:Fefferman}
Charles L. Fefferman. \emph{Existence and Smoothness of the Navier-Stokes Equation}. Clay Mathematics Institute, Millennium Prize Problem statement, \texttt{https:/\allowbreak{}/\allowbreak{}www.claymath.org/\allowbreak{}wp-content/\allowbreak{}uploads/\allowbreak{}2022/\allowbreak{}06/\allowbreak{}navierstokes.pdf}.\newline
{\footnotesize Copy held \texttt{sources/\allowbreak{}mathematics/\allowbreak{}millennium\_\allowbreak{}problems/\allowbreak{}clay\_\allowbreak{}fefferman\_\allowbreak{}navier\_\allowbreak{}stokes.pdf}, SHA-256 \texttt{c1b5f27b1a64705cfaf1afceea513db5deedca8a18ca56ab32e7f86445a06d0c}.}

\bibitem{src:Formal-Conjectures}
Google DeepMind Formal Conjectures authors. \emph{Lean formalization of the Navier-Stokes existence and smoothness problem statement}. \texttt{https:/\allowbreak{}/\allowbreak{}github.com/\allowbreak{}google-deepmind/\allowbreak{}formal-conjectures/\allowbreak{}blob/\allowbreak{}main/\allowbreak{}FormalConjectures/\allowbreak{}Millenium/\allowbreak{}NavierStokes.lean}.

\bibitem{src:Jaffe-and-Witten}
Arthur Jaffe, Edward Witten. \emph{Quantum Yang-Mills Theory}. Clay Mathematics Institute, Millennium Prize Problem statement, \texttt{https:/\allowbreak{}/\allowbreak{}www.claymath.org/\allowbreak{}wp-content/\allowbreak{}uploads/\allowbreak{}2022/\allowbreak{}06/\allowbreak{}yangmills.pdf}.\newline
{\footnotesize Copy held \texttt{sources/\allowbreak{}mathematics/\allowbreak{}millennium\_\allowbreak{}problems/\allowbreak{}clay\_\allowbreak{}jaffe\_\allowbreak{}witten\_\allowbreak{}quantum\_\allowbreak{}yang\_\allowbreak{}mills.pdf}, SHA-256 \texttt{3558403ca14c11e382f73a09e548222708540bfdf478cf96aa11c52d43e23e09}.}

\bibitem{src:Leray}
Leray. 1934.

\bibitem{src:Moore}
Moore. Owed a full citation.

\bibitem{src:OpenAI-2026-Euler}
OpenAI. \emph{Finite time blowup for the Euler equation}. \texttt{https:/\allowbreak{}/\allowbreak{}cdn.openai.com/\allowbreak{}pdf/\allowbreak{}315b36cd-ec98-4023-8342-93345194ece1/\allowbreak{}euler.pdf}. 2026.\newline
{\footnotesize Copy held \texttt{sources/\allowbreak{}mathematics/\allowbreak{}fluid\_\allowbreak{}dynamics/\allowbreak{}openai\_\allowbreak{}2026\_\allowbreak{}finite\_\allowbreak{}time\_\allowbreak{}blowup\_\allowbreak{}euler.pdf}, SHA-256 \texttt{a0c234518e6c489e16996805023eb2e75c00b7c03455f7a3a5be2c124954bfdd}.}

\bibitem{src:OpenAI-2026-Navier-Stokes}
OpenAI. \emph{Finite time blowup for Navier-Stokes}. \texttt{https:/\allowbreak{}/\allowbreak{}cdn.openai.com/\allowbreak{}pdf/\allowbreak{}32d9f210-8b73-45e0-91bc-82a30aef8a9a/\allowbreak{}navier-stokes.pdf}. 2026.\newline
{\footnotesize Copy held \texttt{sources/\allowbreak{}mathematics/\allowbreak{}fluid\_\allowbreak{}dynamics/\allowbreak{}openai\_\allowbreak{}2026\_\allowbreak{}finite\_\allowbreak{}time\_\allowbreak{}blowup\_\allowbreak{}navier\_\allowbreak{}stokes.pdf}, SHA-256 \texttt{0e779481c4da40bd28d1e642e1d8ca57447d129610df28dfa5a11e9af8ae228f}.}

\bibitem{src:Riemann-1859}
Bernhard Riemann. \emph{Handwritten working manuscript, Nachlass Cod. Ms. Riemann 3, Rechenseiten}. Niedersaechsische Staats- und Universitaetsbibliothek Goettingen, Cod. Ms. Riemann 3; scan page 00000038.jpg. 1859.\newline
{\footnotesize Copy held \texttt{sources/\allowbreak{}mathematics/\allowbreak{}analytic\_\allowbreak{}number\_\allowbreak{}theory/\allowbreak{}riemann\_\allowbreak{}nachlass\_\allowbreak{}cod\_\allowbreak{}ms\_\allowbreak{}3\_\allowbreak{}manuscript\_\allowbreak{}scan.pdf}, SHA-256 \texttt{2dd0ea0dde9df70cf6ae441247b1a0716eeaa434a55f9f006d70b48ae10134c2}.}

\bibitem{src:Wiles}
Andrew Wiles. \emph{The Birch and Swinnerton-Dyer Conjecture}. Clay Mathematics Institute, Millennium Prize Problem statement, \texttt{https:/\allowbreak{}/\allowbreak{}www.claymath.org/\allowbreak{}wp-content/\allowbreak{}uploads/\allowbreak{}2022/\allowbreak{}05/\allowbreak{}birchswin.pdf}.\newline
{\footnotesize Copy held \texttt{sources/\allowbreak{}mathematics/\allowbreak{}millennium\_\allowbreak{}problems/\allowbreak{}clay\_\allowbreak{}wiles\_\allowbreak{}birch\_\allowbreak{}swinnerton\_\allowbreak{}dyer.pdf}, SHA-256 \texttt{c25dc966dd0051a60ac04593aef344d8aa89a51abedccbedadc5dc42099ddc6a}.}

\bibitem{src:Wilkins-1998}
Bernhard Riemann, translated by David R. Wilkins. \emph{On the Number of Prime Numbers less than a Given Quantity}. Preliminary version, December 1998, of Ueber die Anzahl der Primzahlen unter einer gegebenen Grosse, Monatsberichte der Berliner Akademie, November 1859.\newline
{\footnotesize Copy held \texttt{sources/\allowbreak{}mathematics/\allowbreak{}analytic\_\allowbreak{}number\_\allowbreak{}theory/\allowbreak{}riemann\_\allowbreak{}1859\_\allowbreak{}wilkins\_\allowbreak{}translation.pdf}, SHA-256 \texttt{6b24341e3406d165baeb7ab1e00cfa0a544d0399e2b33175dfa71480bdddd24f}.}

\end{thebibliography}
